\uuid{DLkV}
\titre{Image directe}
\auteur{bodin ; hanani ; mzari}
\organisation{exo7}
\licence{CC BY-NC-SA 4.0}
\datecreate{2020-01-20}
\niveau{L1}
\module{Algèbre}
\chapitre{Logique, ensemble, raisonnement}
\sousChapitre{Ensemble}
\difficulte{2}
\theme{Ensembles, applications, image directe}
\exo7qcmid{lille-1.ensembles.018}

\contenu{
\begin{question}
\qclassification{100.02,101.01,102.02,102.02}
\qtags{motcle=image directe}
On considère l'application $f:\Rr\to \Rr$ définie par
$$\forall x\in \Rr,\; f(x)=\cos (\pi x).$$
Quelles sont les bonnes réponses ?
\begin{answers}
    \good{$f(\{0,2\})=\{1\}$}
    \bad{$f(\{0,2\})=\{0\}$}
    \bad{$f([0,2])=[1,1]$}
    \good{$f([0,2])=[-1,1]$}
\end{answers}
\begin{explanations}
D'abord, $f(0)=f(2)=1$. Mais, $f$ est décroissante sur $[0,1]$ et est croissante sur $[1,2]$ avec $f(1)=-1$. Dessiner le graphe de $f$ !
\end{explanations}
\end{question}
}
